\section{Workspace Execution}
\label{sec:workspace-execution}

Workspace creation begins by acquiring a snapshot lease, not by reading an
unprotected copy of the current head.  The lease records the selected manifest
and its resolved ordered layer paths; the workspace handle carries that
snapshot together with session-unique runtime directories.  Only after this
view has been fixed does the runtime open the executable workspace.  A later
publication may advance the active manifest, but it does not change the
manifest or layer ordering observed through the live session's lease.  This is
the concrete mechanism behind the stable-view invariant in
Section~\ref{sec:system-model}, without importing serializable snapshot
isolation.

On the supported Linux path, the runtime projects the leased LayerStack with
OverlayFS.  The manifest's layers become read-only lower directories in
newest-first precedence order, while the session receives a unique upper
directory and work directory.  The resulting mount therefore reads through
shared history but records regular writes, copied-up files, and deletion or
opacity metadata only in the private upper tree.  Lower layers are shared
runtime state; the upper/work pair is session-private state.  This mechanism
is Linux- and OverlayFS-centered at the audited baseline: the non-Linux mount
path reports the operation as unsupported, and the paper makes no universal
cross-platform execution claim.

Each workspace also has a long-lived namespace-holder process.  The holder
creates and keeps open the session's user, mount, and PID namespaces and, when
selected, its network namespace.  The daemon retains file descriptors for
these namespaces rather than moving its own control process into them.  For
each command, a short-lived runner joins the recorded user, mount, PID, and
optional network namespaces in the runtime's required order and launches the
command against the already mounted workspace.  This holder/runner split lets
successive commands enter the same filesystem and process-namespace context
while the daemon remains outside the command namespaces.  It is an
implementation isolation boundary, not a proof that hostile code cannot
escape or interfere through every host resource.

A sessionless \texttt{exec\_command} uses the \emph{implicit session} path.
The operation service creates a workspace with the shared-network profile and
a publish-then-destroy finalization policy, admits the requested command, and
tracks it in the session's command ledger.  Publication is eligible only after
the active-command set becomes empty and the ledger reaches its drained
terminal state; it is therefore tied to command completion, not merely to
returning an initial command handle or response.  The finalizer then attempts
the capture, publication, and destruction phases described below.  This
automatic lifecycle is specific to sessionless command execution.

An \emph{explicit session} instead keeps one leased view, namespace holder,
mount, and private upper/work pair across multiple calls.  Commands admitted
with its session identifier and file operations directed to that session
observe the same live overlay, so later operations can consume earlier private
writes before anything reaches the active head.  The explicit session uses a
no-op automatic finalization policy: its caller deliberately requests
publication or destruction.  Publication admission fails while commands are
still active, preventing capture from racing an admitted command through the
normal service path.

File and network qualifiers matter to this account.  A session-scoped read,
write, or edit operates on the live mounted workspace and therefore on the
private overlay.  A sessionless read instead projects the latest published
LayerStack, while a sessionless write or edit directly amends the current head
with a one-layer operation under the LayerStack writer lock; neither
sessionless file path creates an implicit workspace session.  Likewise,
shared networking joins the host network namespace while retaining separate
user, mount, and PID namespaces, whereas the isolated profile creates a
separate network namespace and its configured connectivity.  Shared
networking is consequently not an egress-containment guarantee, and selectable
network isolation is not evidence of universal denial policy.

Workspace creation has several externally fallible steps---lease acquisition,
directory allocation, holder startup, namespace access, overlay mounting, and
optional network setup.  The service attempts to roll back resources created
before a failed admission, while a successfully admitted session retains its
lease and private tree until its lifecycle reaches cleanup.  Neither path
captures command registers, memory, sockets, or an arbitrary process tree for
later rollback.  Publication begins only from the filesystem delta left in the
complete private upper directory; the next section follows that delta from
capture to the active head.
